\subsection{Tests and complete-call measurements}

All 38 focused tests pass in 4.469 seconds, and all 1,214 maintained tests
pass in 235.344 seconds. The final source audit takes 135.085 seconds and
the cell-count survey takes 20.777 seconds. Those correctness runs partly
overlap the test suite; their elapsed times are not paired performance
measurements. The audit, survey and both timing runs pin 496 source files.
The six existing seed, cocycle, annulus, planar, span and weighted-orbit
verifier modules remain byte-identical to the baseline revision.

Timing endpoints start from input diagrams and include complete
recognition, any exact fallback, independent replay of returned native
certificates and certificate hashing. Earlier shortcuts are disabled to
reach the native stage. Both implementations have no local seed cap and
retain the twelve-second global recognition deadline. The comparison
selects the published or current native search body while retaining the
same public construction and verification APIs. No tests, audits, builds
or commits overlap the timed runs.

The first five-round experiment shows noticeable variation in the raw-disc
old-arm A/A ratio, 1.115, and one later-tree control, 1.074. It is retained,
and a nine-round repeat supplies the final table. Each experiment includes
one warm-up round per case and randomizes old, old A/A, new and new A/A
calls within each round. All 616 measured calls and 88 warm-ups complete.
Per-arm verdicts, certificate hashes and native work counts are stable
across both runs. The initial run takes 24.056 seconds; the final run,
with 396 measured calls and 44 warm-ups, takes 49.367 seconds.

The two later-tree annuli and the later planar cap request 24 tree trials;
other cases request four. Only the two labeled planar cases enable planar
queries. The face-disc case disables annulus recognition and requests four
face roots, so it exercises the actual face-search path. All other cases
retain the existing annulus criterion and zero face roots. The table shows
primary-arm median times and median paired ratios; these ratios need not
equal quotients of the displayed median times.

\begin{center}
\small\begin{tabular}{lrrrrr}
\toprule
Input & Old (ms) & New (ms) & Old/new & Old A/A & New A/A \\
\midrule
Immediate disc & 62.879 & 64.832 & 0.983 & 1.000 & 0.988 \\
Immediate annulus & 31.083 & 31.531 & 0.980 & 1.000 & 1.004 \\
Optimized disc & 90.827 & 70.859 & 1.281 & 1.002 & 1.002 \\
Optimized annulus & 136.156 & 108.234 & 1.262 & 1.001 & 1.006 \\
Later tree annulus A & 105.399 & 67.743 & 1.509 & 0.998 & 0.999 \\
Later tree annulus B & 88.082 & 60.483 & 1.466 & 1.012 & 1.000 \\
Early planar cap & 91.193 & 73.902 & 1.283 & 1.011 & 1.035 \\
Later planar cap & 426.389 & 362.727 & 1.179 & 1.021 & 1.020 \\
Optional face disc & 222.422 & 152.559 & 1.400 & 0.990 & 1.002 \\
Trefoil miss & 42.453 & 20.975 & 2.002 & 0.996 & 0.998 \\
Genus-one miss & 43.638 & 23.442 & 1.887 & 1.010 & 1.005 \\
\bottomrule
\end{tabular}

\end{center}

The optimized disc and annulus improve by paired factors 1.281 and 1.262.
The two later-tree annuli improve 1.509 and 1.466. Optional planar cases
improve 1.283 and 1.179, and the optional face-disc endpoint improves 1.400.
The trefoil and genus-one continuing misses improve 2.002 and 1.887,
including the same exact fallback used before. These are gains on complete
calls through the selected native paths, rather than just on a standalone
geometry helper.

The immediate disc and annulus are adverse controls: their paired ratios
are 0.983 and 0.980, approximately two percent slower. They pay for the
prefilter before their existing independent replay and avoid none of the
old negative replay work. There is no claim of a uniform speedup. Absolute
late-planar times also vary substantially between runs, while its paired
factor remains about 1.18. Comparing raw elapsed levels across experiments
would overstate what those measurements establish.

The default candidate schedule is retained because discovery geometry
reuse improves continuing searches while preserving every positive gate;
no additional search is introduced. The planar and face queries remain
opt-in, and the whole native stage remains optional in the ordinary
portfolio. Its nineteen control sources still finish before that stage,
so this experiment supplies no new ordinary-portfolio speedup claim.

Reproduce using \path{../fast/normal_orbit_research/reuse.py} in
\code{audit}, \code{survey} or \code{recognize --rounds 9} mode. The
source audit verifies that the loaded baseline search reproduces the
previous records before comparing the new implementation. The checked
incoming inventory and \code{openai/math} head \code{fd4aeeb2e} are
unchanged. This removes redundant polynomial work and supplies a linear
cell-count prefilter; it does not complete the surface candidate search,
normal capped-disc extraction or hierarchy operations, and establishes no
general subexponential or quasipolynomial recognition bound.
